\section{Definitions}
A complex number $z$ is generally written in the form $z=a+bi$ where $a$ and $b$ are real numbers and $i$, called the imaginary unit, has the property that $i^2=-1$. The real numbers $a$ and $b$ are called the real and imaginary parts of $z=a+bi$, respectively.
The complex conjugate of $z$ is denoted by $\bar{z}$; it is defined by $\overline{a+bi}=a-bi$. Thus, $a+bi$ and $a-bi$ are conjugates of each other.
\section{Equality of Complex Numbers}
$a+bi=c+di$ if and only if $a=c$ and $b=d$
\section{Arithmetic of Complex Numbers}
$$(a+bi)+(c+di)=(a+c)+(b+d)i$$
$$(a+bi)-(c+di)=(a-c)+(b-d)i$$
$$(a+bi)(c+di)=(ac-bd)+(ad+bc)i$$
$$\dfrac{a+bi}{c+di}=\dfrac{(a+bi)(c-di)}{(c+di)(c-di)}=\dfrac{ac+bd}{c^2+d^2}+\dfrac{bc-ad}{c^2+d^2}i$$
